\chapter*{前言}
\label{cha:preface}

% Uncomment this if you think Preface should appear in TOC
%\addcontentsline{toc}{chapter}{前言}
\subsection*{普林斯顿高等研究院 泛等基础特别学年}

2012--13 学年，在 Steve Awodey、Thierry Coquand 和 Vladimir Voevodsky 的组织下，普林斯顿高等研究院数学学院举办了一场关于数学的泛等基础的特别学年活动。以下是官方参与者名单。

\begin{multicols}{\OPTprefacecols}{
\begin{itemize}
\item[] Peter Aczel
\item[] Benedikt Ahrens
\item[] Thorsten Altenkirch
\item[] Steve Awodey
\item[] Bruno Barras
\item[] Andrej Bauer
\item[] Yves Bertot
\item[] Marc Bezem
\item[] Thierry Coquand
\item[] Eric Finster
\item[] Daniel Grayson
\item[] Hugo Herbelin
\item[] Andr\'e Joyal
\item[] Dan Licata
\item[] Peter Lumsdaine
\item[] Assia Mahboubi
\item[] Per Martin-L\"of
\item[] Sergey Melikhov
\item[] Alvaro Pelayo
\item[] Andrew Polonsky
\item[] Michael Shulman
\item[] Matthieu Sozeau
\item[] Bas Spitters
\item[] Benno van den Berg
\item[] Vladimir Voevodsky
\item[] Michael Warren
\item[] Noam Zeilberger
\end{itemize}
}
\end{multicols}

\noindent 此外，还有以下学生参与，他们的贡献同样宝贵。

\begin{multicols}{\OPTprefacecols}{
\begin{itemize}
\item[] Carlo Angiuli
\item[] Anthony Bordg
\item[] Guillaume Brunerie
\item[] Chris Kapulkin
\item[] Egbert Rijke
\item[] Kristina Sojakova
\end{itemize}
}
\end{multicols}

\noindent 另外，还有以下短期与长期访问学者（包括学生访客），他们对特别学年的贡献同样至关重要。

\begin{multicols}{\OPTprefacecols}{
\begin{itemize}
\item[] Jeremy Avigad
\item[] Cyril Cohen
\item[] Robert Constable
\item[] Pierre-Louis Curien
\item[] Peter Dybjer
\item[] Mart{\'\i}n Escard{\'o}
\item[] Kuen-Bang Hou
\item[] Nicola Gambino
\item[] Richard Garner
\item[] Georges Gonthier
\item[] Thomas Hales
\item[] Robert Harper
\item[] Martin Hofmann
\item[] Pieter Hofstra
\item[] Joachim Kock
\item[] Nicolai Kraus
\item[] Nuo Li
\item[] Zhaohui Luo
\item[] Michael Nahas
\item[] Erik Palmgren
\item[] Emily Riehl
\item[] Dana Scott
\item[] Philip Scott
\item[] Sergei Soloviev
\end{itemize}
}
\end{multicols}

\subsection*{关于本书}

我们起初并没有打算写一本书。
本书源于大家的一番共同努力：提出一种新的、人类可读可理解“非形式化类型论”写法，作为可由机器验证的形式证明的补充。
泛等基础与这样一种数学基础观紧密相连：数学应该基于一套能够在计算机证明助理中实现的基础。
虽然本书本身并不包含这样的形式化工作，但这里介绍的许多内容，实际上最初都是在证明助理中以完全形式化的方式完成的；
直到后来，我们才将其“去形式化”，整理成如今读者眼前的模样。这恰好将数学形式化中的通常次序颠倒了过来，实属一件奇事。

前述的每一位参与者，都以自己的方式为那一“特别学年”作出了贡献——也因此为本书作出了贡献；这些贡献或是思想，或是文字，或是实际的工作。
贯穿整整一年的那种合作精神，确实非同寻常。

\mentalpause

在此要特别感谢高等研究院；没有它，这本书显然根本不可能问世。
事实证明，这里正是孕育这个数学新分支的理想环境：思想活跃，气氛融洽，鼎力支持。
但愿这种独一无二的氛围还能有一丝留存于本书的字里行间，也留存在这一新兴领域的未来发展之中。

\bigskip

\begin{flushright}
The Univalent Foundations Program\\
普林斯顿高等研究院\\
普林斯顿，2013年4月
\end{flushright}

%%%%%%%%%%%%% end of scope of local macros
% Local Variables:
% TeX-master: "hott-online"
% End: